
This work attempted to develop efficient semantic entropy proxies via similarity-augmented hidden state probes. The experiment produced negative results: the tested configuration (layer 25, TBG position, logistic regression, Llama-3-8B-Instruct) achieved Spearman correlation rho = 0.0843 with true SE, statistically indistinguishable from zero, and AUROC = 0.52, indistinguishable from random classification.

The 39\% AUROC gap between these results and published benchmarks indicates substantial configuration sensitivity. The experiment tested one configuration without systematic ablation; the negative result applies to this specific setup rather than to SE probing in general.

The contributions of this work are:

\begin{enumerate}
\item Documentation of SE probe failure under a specific configuration, demonstrating that this configuration does not yield meaningful SE prediction on TruthfulQA.
\end{enumerate}

\begin{enumerate}
\item Identification of a replication gap that highlights potential configuration sensitivity in SE probing.
\end{enumerate}

\begin{enumerate}
\item A negative result that may inform future work by emphasizing the need for systematic configuration validation.
\end{enumerate}

The failure of this configuration does not preclude success under alternative configurations. Systematic layer ablation, token position comparison, and probe architecture exploration remain as future work that may identify configurations where SE probing succeeds.
